\documentclass[../../main.tex]{subfiles}% 注意这里的文件路径不能用 ./main.tex ,否则用latexmk编译子文件会报错
\graphicspath{{\subfix{./image/}}{\subfix{../../image/}}} % 指定图片引用路径,后续可以直接使用图片文件名
% 如果图片存放在上级目录,则使用 ../../image/ 作为备用路径

% 例如:
% \begin{figure}[H]
% \centering
% \includegraphics[width=0.3\textwidth]{图.png}
% \caption{}
% \label{figure:图}
% \end{figure}
% 注意:上述\label{}一定要放在\caption{}之后,否则引用图片序号会只会显示??.

\begin{document}

\section{习题}

\begin{example}
用特征线法求解下列 Cauchy 初值问题：
\begin{enumerate}[(1)]
\item $\begin{cases}
u_t+2u_x=0, & (x,t)\in\mathbb{R}_+^2, \\
u(x,0)=x^2, & x\in\mathbb{R};
\end{cases}$
\item $\begin{cases}
u_t+2u_x+u=xt, & (x,t)\in\mathbb{R}_+^2, \\
u(x,0)=2-x, & x\in\mathbb{R};
\end{cases}$
\item $\begin{cases}
2u_t-u_x+xu=0, & (x,t)\in\mathbb{R}_+^2, \\
u(x,0)=2xe^{x^2/2}, & x\in\mathbb{R};
\end{cases}$
\item $\begin{cases}
u_t+(1+x^2)u_x-u=0, & (x,t)\in\mathbb{R}_+^2, \\
u(x,0)=\arctan x, & x\in\mathbb{R};
\end{cases}$
\item $\begin{cases}
u_t+A\cdot Du+cu=0, & (x,t)\in\mathbb{R}_+^{n+1}, \\
u(x,0)=\varphi(x), & x\in\mathbb{R}^n,
\end{cases}$
这里 $A\in\mathbb{R}^n$ 是常向量，$c\in\mathbb{R}$ 是常数。
\end{enumerate}
\end{example}

\begin{example}
用特征线法求解拟线性方程的 Cauchy 初值问题
\begin{align*}
\begin{cases}
u_t+uu_x=0, & (x,t)\in\mathbb{R}_+^2, \\
u(x,0)=\varphi(x), & x\in\mathbb{R},
\end{cases}
\end{align*}
其中 $\varphi(x)$ 是一个递增函数。（提示：考虑特征线 $\frac{dx}{dt}=u$）
\end{example}

\begin{example}
\begin{enumerate}[(1)]
\item 设 $u=u(x,y)$ 在凸连通区域 $\Omega\subset\mathbb{R}^2$ 上满足方程 $u_{xy}=0$，求此方程的所有解。
\item 设 $a$ 为正常数。利用变量替换 $\xi=x+at,\ \eta=x-at$ 证明 $u_{tt}-a^2u_{xx}=0$ 当且仅当 $u_{\xi\eta}=0$。
\item 利用 (1), (2) 推导 D'Alembert 公式。
\end{enumerate}
\end{example}

\begin{example}
设 $ABCD$ 表示 $\mathbb{R}_+^2$ 上的一个平行四边形，其中两条边 $AB,CD$ 平行于波动方程 $u_{tt}-a^2u_{xx}=0$ ($a>0$) 的一条特征线 $x-at=0$，另两条边 $AD,BC$ 平行于另一条特征线 $x+at=0$。这样的平行四边形称为特征平行四边形。设 $u(x,t)$ 是波动方程的解。证明：
\begin{align*}
u(A)+u(C)=u(B)+u(D),
\end{align*}
其中 $u(A),u(B),u(C),u(D)$ 分别表示函数 $u$ 在顶点 $A,B,C,D$ 处的函数值。
\end{example}

\begin{example}
假设电场强度 $\bm{E}=(E_1,E_2,E_3)$ 和磁场强度 $\bm{B}=(B_1,B_2,B_3)$ 满足 Maxwell 方程组
\begin{align*}
\begin{cases}
\frac{1}{c}\cdot\frac{\partial \bm{E}}{\partial t}=\operatorname{curl}\bm{B}, \\
\frac{1}{c}\cdot\frac{\partial \bm{B}}{\partial t}=-\operatorname{curl}\bm{E}, \\
\operatorname{div}\bm{E}=\operatorname{div}\bm{B}=0,
\end{cases}
\end{align*}
这里 $c$ 为光速。证明：它们的分量 $E_1,E_2,E_3,B_1,B_2,B_3$ 满足波动方程 $u_{tt}-c^2\Delta u=0$。
\end{example}

\begin{example}
证明：方程
\begin{align*}
\left(1-\frac{x}{h}\right)^2\frac{\partial^2u}{\partial t^2}
=a^2\frac{\partial}{\partial x}\left[\left(1-\frac{x}{h}\right)^2\frac{\partial u}{\partial x}\right]
\end{align*}
的通解可以表示成
\begin{align*}
u(x,t)=\frac{F(x-at)+G(x+at)}{h-x},
\end{align*}
其中 $F,G$ 是任意二次连续可微函数，$h>0,\ a>0$ 为常数。（提示：考虑函数替换 $v=(h-x)u$）
\end{example}

\begin{example}
直接用特征线法来求解初值问题 \eqref{eq:wave-cauchy1d}。
\end{example}

\begin{example}
直接验证定理 \eqref{eq:wave-cauchy1d} 和推论。
\end{example}

\begin{example}
当初值 $u(x,0)=\varphi(x),\ u_t(x,0)=\psi(x)$ 满足什么条件时，一维齐次波动方程的初值问题的解仅由右行波组成（即通解为 $G(x-at)$ 的形式）？
\end{example}

\begin{example}
已知初值问题
\begin{align*}
\begin{cases}
u_{tt}-u_{xx}=0, & x\in\mathbb{R},\ t>a, \\
u|_{t=a}=u_0(x), & x\in\mathbb{R}, \\
u_t|_{t=a}=u_1(x), & x\in\mathbb{R},
\end{cases}
\end{align*}
其中 $a\neq\pm 1$。若初值在 $b\leqslant x\leqslant c$ 上给定，试问它的解在什么区域确定。
\end{example}

\begin{example}
证明：初值问题
\begin{align*}
\begin{cases}
u_{tt}-u_{xx}=6(x+t), & x\in\mathbb{R},\ t>x, \\
u|_{t=x}=0, & x\in\mathbb{R}, \\
u_t|_{t=x}=\psi(x), & x\in\mathbb{R}
\end{cases}
\end{align*}
有解的充分必要条件是 $\psi(x)-3x^2=\text{const}$；若上述问题有解，则问题的解不是唯一的。试问：当在直线 $t=ax$ 上给定初值时，为什么在 $a=\pm 1$ 和 $a\neq\pm 1$ 的情形，关于解的存在性唯一性的结论完全不同？
\end{example}

\begin{example}
设 $a>0$ 为常数。利用 Fourier 变换求解三维波动方程初值问题
\begin{align*}
\begin{cases}
u_{tt}-a^2(u_{xx}+u_{yy}+u_{zz})=f(x,y,z,t), & (x,y,z,t)\in\mathbb{R}^4, \\
u|_{t=0}=\varphi(x,y,z), & (x,y,z)\in\mathbb{R}^3, \\
u_t|_{t=0}=\psi(x,y,z), & (x,y,z)\in\mathbb{R}^3.
\end{cases}
\end{align*}
\end{example}

\begin{example}
设 $a>0$ 为常数。利用 Fourier 变换求解二维波动方程初值问题
\begin{align*}
\begin{cases}
u_{tt}-a^2(u_{xx}+u_{yy})=f(x,y,t), & (x,y,t)\in\mathbb{R}_+^3, \\
u|_{t=0}=\varphi(x,y), & (x,y)\in\mathbb{R}^2, \\
u_t|_{t=0}=\psi(x,y), & (x,y)\in\mathbb{R}^2.
\end{cases}
\end{align*}
\end{example}

\begin{example}
设 $u=u(x,y,z,t)$ 是三维波动方程初值问题
\begin{align*}
\begin{cases}
u_{tt}-a^2(u_{xx}+u_{yy}+u_{zz})=0, & (x,y,z,t)\in\mathbb{R}_+^4, \\
u|_{t=0}=f(x)+g(y), & (x,y,z)\in\mathbb{R}^3, \\
u_t|_{t=0}=\varphi(y)+\psi(z), & (x,y,z)\in\mathbb{R}^3
\end{cases}
\end{align*}
的解，其中 $a>0$ 为常数。求解 $u$ 的表达式。
\end{example}

\begin{example}
设函数 $\Phi\in C^2(\mathbb{R})$，向量 $\alpha\in\mathbb{R}^n$ 且 $|\alpha|=1$，则 $\Phi(\alpha\cdot x+at)$ 满足 $n$ 维波动方程
\begin{align*}
u_{tt}-a^2\Delta u=0,
\end{align*}
这里 $a>0$ 为常数，$x=(x_1,\cdots,x_n)$，$\Delta=\frac{\partial^2}{\partial x_1^2}+\cdots+\frac{\partial^2}{\partial x_n^2}$。波动方程的这种形式的解称为平面波解。
\end{example}

\begin{example}
设 $a$ 为正常数。证明：三维波动方程
\begin{align*}
u_{tt}-a^2\Delta u=0,\quad (x,t)\in\mathbb{R}_+^4
\end{align*}
的所有径向对称的解可以表示为
\begin{align*}
u(x,t)=\frac{F(|x|-at)+G(|x|+at)}{|x|},\quad x\neq 0.
\end{align*}
\end{example}

\begin{example}
设 $a$ 为正常数。求解下列多维初值问题：
\begin{enumerate}[(1)]
\item $\begin{cases}
u_{tt}-a^2(u_{xx}+u_{yy})=0, & (x,y,t)\in\mathbb{R}_+^3, \\
u|_{t=0}=x^2(x+y),\quad u_t|_{t=0}=0, & (x,y)\in\mathbb{R}^2;
\end{cases}$
\item $\begin{cases}
u_{tt}-a^2(u_{xx}+u_{yy}+u_{zz})=0, & (x,y,z,t)\in\mathbb{R}_+^4, \\
u|_{t=0}=x^2+y^2z,\quad u_t|_{t=0}=1+y, & (x,y,z)\in\mathbb{R}^3.
\end{cases}$
\end{enumerate}
\end{example}

\begin{example}
设 $a$ 为正常数，函数 $\varphi(x),\psi(x)\in C^2[0,+\infty)$，$g(t)\in C^2[0,+\infty)$ 满足相容性条件
\begin{align*}
g(0)=\varphi(0),\quad g'(0)=\psi(0),\quad g''(0)=a^2\varphi''(0).
\end{align*}
试求解半无界问题
\begin{align*}
\begin{cases}
u_{tt}-a^2u_{xx}=0, & (x,t)\in\mathbb{R}_+\times\mathbb{R}_+, \\
u(x,0)=\varphi(x),\quad u_t(x,0)=\psi(x), & x\geqslant 0, \\
u(0,t)=g(t), & t\geqslant 0.
\end{cases}
\end{align*}
\end{example}

\begin{example}
试给出半无界问题
\begin{align*}
\begin{cases}
u_{tt}-a^2u_{xx}=f(x,t), & (x,t)\in\mathbb{R}_+\times\mathbb{R}_+, \\
u(x,0)=\varphi(x),\quad u_t(x,0)=\psi(x), & x\geqslant 0, \\
u_x(0,t)=g(t), & t\geqslant 0
\end{cases}
\end{align*}
有古典解的相容性条件。
\end{example}

\begin{example}
利用惟一性结果直接证明：当初值 $\varphi(x),\psi(x)$ 和非齐次项 $f(x,t)$ 是关于 $x$ 的偶函数时初值问题 \eqref{eq:wave-cauchy1d} 的解 $u(x,t)$ 也是关于 $x$ 的偶函数。根据以上事实，用延拓法求解半无界问题
\begin{align*}
\begin{cases}
u_{tt}-a^2u_{xx}=f(x,t), & (x,t)\in\mathbb{R}_+\times\mathbb{R}_+, \\
u(x,0)=u_t(x,0)=0, & x\geqslant 0, \\
u_x(0,t)=0, & t\geqslant 0,
\end{cases}
\end{align*}
并说明当 $f(x,t)$ 满足什么条件时，导出的形式解确实是问题的解。这里 $a$ 为正常数。
\end{example}

\begin{example}
利用延拓法求解半无界问题
\begin{align*}
\begin{cases}
u_{tt}-a^2u_{xx}=0, & (x,t)\in\mathbb{R}_+\times\mathbb{R}_+, \\
u(x,0)=u_t(x,0)=0, & x\geqslant 0, \\
u_x(0,t)=A\sin\omega t, & t\geqslant 0,
\end{cases}
\end{align*}
并给出物理解释。这里 $a$ 为正常数，$A$ 为常数。
\end{example}

\begin{example}
设 $u=u(x,y,t)$ 是初值问题
\begin{align*}
\begin{cases}
u_{tt}-4(u_{xx}+u_{yy})=0, & (x,y,t)\in\mathbb{R}_+^3, \\
u(x,y,0)=0, & (x,y)\in\mathbb{R}^2, \\
u_t(x,y,0)=\psi(x,y), & (x,y)\in\mathbb{R}^2
\end{cases}
\end{align*}
的解，其中当 $(x,y)\in\Omega$ 时，$\psi(x,y)\equiv 0$；当 $(x,y)\in\mathbb{R}^2\setminus\Omega$ 时，$\psi(x,y)>0$。这里 $\Omega$ 是正方形 $\{(x,y)\in\mathbb{R}^2: |x|\leqslant 1,\ |y|\leqslant 1\}$。试指出当 $t>0$ 时，$u(x,y,t)\equiv 0$ 的区域。
\end{example}

\begin{example}
设 $u=u(x,t)$ 是半无界问题
\begin{align*}
\begin{cases}
u_{tt}-a^2u_{xx}=0, & (x,t)\in\mathbb{R}_+\times\mathbb{R}_+, \\
u(x,0)=\varphi(x), & x\geqslant 0, \\
u_t(x,0)=\psi(x), & x\geqslant 0, \\
u(0,t)=g(t), & t\geqslant 0
\end{cases}
\end{align*}
的解，其中 $a$ 为正常数，且
\begin{align*}
\varphi(x),\psi(x)=\begin{cases}
0, & x\in[0,1], \\
>0, & x\in(1,+\infty),
\end{cases}\qquad
g(t)=\begin{cases}
0, & t\in[0,1], \\
>0, & t\in(1,+\infty).
\end{cases}
\end{align*}
试指出当 $t>0$ 时，$u(x,t)\equiv 0$ 的区域。
\end{example}

\begin{example}
试问：半无界问题
\begin{align*}
\begin{cases}
u_{tt}-u_{xx}+u_t+u_x=0, & (x,t)\in\mathbb{R}_+\times\mathbb{R}_+, \\
u(x,0)=\varphi(x), & x\geqslant 0, \\
u_t(x,0)=\psi(x), & x\geqslant 0, \\
u(0,t)=0, & t\geqslant 0
\end{cases}
\end{align*}
能否直接用对称开拓法求解？为什么？试用特征线法求解此半无界问题。
提示：利用等式
\begin{align*}
\frac{\partial^2}{\partial t^2}-\frac{\partial^2}{\partial x^2}+\frac{\partial}{\partial t}+\frac{\partial}{\partial x}
=\left(\frac{\partial}{\partial t}+\frac{\partial}{\partial x}\right)\left(\frac{\partial}{\partial t}-\frac{\partial}{\partial x}+1\right)
\end{align*}
将方程化为两个一阶的方程。
\end{example}

\begin{example}
求解 Goursat 问题
\begin{align*}
\begin{cases}
u_{tt}-u_{xx}=0, & t>|x|, \\
u|_{t=-x}=\varphi(x), & x\leqslant 0, \\
u|_{t=x}=\psi(x), & x\geqslant 0,
\end{cases}
\end{align*}
其中 $\varphi(0)=\psi(0)$。如果 $\varphi(x)$ 在 $(-a,0]$ ($a>0$) 上给定，$\psi(x)$ 在 $[0,b]$ ($b>0$) 上给定，试指出此定解条件的决定区域。
\end{example}

\begin{example}
求解 Darboux 问题
\begin{align*}
\begin{cases}
u_{tt}-u_{xx}=0, & 0<x<t, \\
u|_{x=0}=\varphi(t), & t\geqslant 0, \\
u|_{x=t}=\psi(t), & t\geqslant 0,
\end{cases}
\end{align*}
其中 $\varphi(0)=\psi(0)$。如果 $\varphi(t),\psi(t)$ 都在 $[0,a]$ ($a>0$) 上给定，试指出此定解条件的决定区域。
\end{example}

\begin{example}
求解问题
\begin{align*}
\begin{cases}
u_{tt}-u_{xx}=0, & 0<x<t, \\
u|_{x=t}=\varphi(t), & t\geqslant 0, \\
u|_{x=0}=\psi(t), & t\geqslant 0.
\end{cases}
\end{align*}
如果 $\varphi(t),\psi(t)$ 都在 $[0,a]$ ($a>0$) 上给定，试指出此定解条件的决定区域。
\end{example}

\begin{example}
求解问题
\begin{align*}
\begin{cases}
u_{tt}-u_{xx}=0, & 0<x<t, \\
u|_{x=t}=\varphi(t), & t\geqslant 0, \\
(-u_x+u)|_{x=0}=\psi(t), & t\geqslant 0.
\end{cases}
\end{align*}
如果 $\varphi(t),\psi(t)$ 都在 $[0,a]$ ($a>0$) 上给定，试指出此定解条件的决定区域。
\end{example}

\begin{example}
在区域 $[0,+\infty)\times[0,+\infty)$ 上考虑一阶方程
\begin{align*}
u_t+au_x=0 \quad\text{或}\quad u_t-au_x=0
\end{align*}
满足初边值条件
\begin{align*}
u|_{x=0}=g(t),\quad t\geqslant 0,\qquad u|_{t=0}=\varphi(x),\quad x\geqslant 0
\end{align*}
的初边值问题，这里 $a>0$ 是常数。试问对于哪一个方程上述初边值条件提法正确？为什么？对提法正确的初边值问题求解，并说明对 $g(t),\varphi(x)$ 提什么条件，所得的解是属于 $C^1$ 的解。
\end{example}

\begin{example}
求解并证明初值问题
\begin{align*}
\begin{cases}
u_{tt}-a^2u_{xx}+bu=f(x,t), & (x,t)\in\mathbb{R}_+^2, \\
u|_{t=0}=\varphi(x), & x\in\mathbb{R}, \\
u_t|_{t=0}=\psi(x), & x\in\mathbb{R}
\end{cases}
\end{align*}
解的惟一性，其中 $a\in\mathbb{R}_+,\ b\in\mathbb{R}$ 为常数。
提示：考虑二维初值问题
\begin{align*}
\begin{cases}
\tilde{u}_{tt}-a^2(\tilde{u}_{xx}+\tilde{u}_{yy})=\tilde{f}(x,y,t), & (x,y,t)\in\mathbb{R}_+^3, \\
\tilde{u}|_{t=0}=\tilde{\varphi}(x,y), & (x,y)\in\mathbb{R}^2, \\
\tilde{u}_t|_{t=0}=\tilde{\psi}(x,y), & (x,y)\in\mathbb{R}^2,
\end{cases}
\end{align*}
其中 $\tilde{u}=\tilde{u}(x,y,t)=u(x,t)v(y)$，$\tilde{f}=f(x,t)v(y)$，$\tilde{\varphi}=\varphi(x)v(y)$，$\tilde{\psi}=\psi(x)v(y)$，辅助函数 $v(y)$ 满足 $v''(y)=-\frac{b}{a^2}v(y)$。
\end{example}

\begin{example}
求解并证明初值问题
\begin{align*}
\begin{cases}
u_{tt}-a^2u_{xx}+bu_t+cu_x+du=f(x,t), & (x,t)\in\mathbb{R}_+^2, \\
u|_{t=0}=\varphi(x), & x\in\mathbb{R}, \\
u_t|_{t=0}=\psi(x), & x\in\mathbb{R}
\end{cases}
\end{align*}
解的惟一性。
\end{example}

\begin{example}
设 $u\in C^2(\mathbb{R}_+^2)$ 是 Cauchy 初值问题
\begin{align*}
\begin{cases}
u_{tt}-a^2u_{xx}=0, & (x,t)\in\mathbb{R}_+^2, \\
u(x,0)=\varphi(x), & x\in\mathbb{R}, \\
u_t(x,0)=\psi(x), & x\in\mathbb{R}
\end{cases}
\end{align*}
的解，其中 $a$ 为正常数，$\varphi\in C^2(\mathbb{R}),\ \psi\in C^1(\mathbb{R})$。证明：
\begin{enumerate}[(1)]
\item 对于任意 $t\geqslant 0$，成立能量不等式
\begin{align*}
\int_{-\infty}^{+\infty}[u_t^2(x,t)+a^2u_x^2(x,t)]\,\mathrm{d}x
\leqslant \int_{-\infty}^{+\infty}[\psi^2(x)+a^2|\varphi'(x)|^2]\,\mathrm{d}x;
\end{align*}
\item 对于任意 $t\geqslant 0$，成立能量恒等式
\begin{align*}
\int_{-\infty}^{+\infty}[u_t^2(x,t)+a^2u_x^2(x,t)]\,\mathrm{d}x
= \int_{-\infty}^{+\infty}[\psi^2(x)+a^2|\varphi'(x)|^2]\,\mathrm{d}x.
\end{align*}
\end{enumerate}
\end{example}

\begin{example}
设函数 $u\in C^2(\mathbb{R}_+^2)$ 满足第 37 题的 Cauchy 初值问题且 $\varphi(x),\psi(x)\in C_0^\infty(\mathbb{R})$。记其动能和势能分别为
\begin{align*}
k(t)=\frac{1}{2}\int_{-\infty}^{+\infty}u_t^2(x,t)\,\mathrm{d}x,\qquad
p(t)=\frac{a^2}{2}\int_{-\infty}^{+\infty}u_x^2(x,t)\,\mathrm{d}x.
\end{align*}
证明：
\begin{enumerate}[(1)]
\item $k(t)+p(t)$ 是与 $t$ 无关的常数；
\item 当 $t$ 足够大时，$k(t)=p(t)$。
\end{enumerate}
提示：利用 D'Alembert 公式和附注 4.2 的表达式，则
\begin{align*}
u_t(x,t)=aF'(x+at)-aG'(x-at),\qquad
u_x(x,t)=F'(x+at)+G'(x-at),
\end{align*}
其中
\begin{align*}
F'=\frac{1}{2a}(a\varphi'+\psi),\qquad G'=\frac{1}{2a}(a\varphi'-\psi).
\end{align*}
\end{example}

\begin{example}
设函数 $u\in C^2(\mathbb{R}_+^4)$ 满足 Cauchy 初值问题
\begin{align*}
\begin{cases}
u_{tt}-a^2\Delta u=0, & (x,t)\in\mathbb{R}_+^4, \\
u(x,0)=\varphi(x), & x\in\mathbb{R}^3, \\
u_t(x,0)=\psi(x), & x\in\mathbb{R}^3,
\end{cases}
\end{align*}
其中 $a$ 为正常数，$\varphi(x),\psi(x)\in C_0^\infty(\mathbb{R}^3)$。证明：存在常数 $C$，使得对于所有的 $(x,t)\in\mathbb{R}_+^4$，成立
\begin{align*}
|u(x,t)|\leqslant \frac{C}{t}.
\end{align*}
（提示：利用 Kirchhoff 公式 \eqref{eq:kirchhoff-hom}）
\end{example}

\begin{example}
设 $a,h$ 为正常数，$A_1,A_2$ 为常数。用分离变量法求解下列混合问题：
\begin{enumerate}[(1)]
\item $\begin{cases}
u_{tt}-a^2u_{xx}=0, & 0<x<\pi,\ t>0, \\
u(x,0)=\sin x,\quad u_t(x,0)=x(\pi-x), & 0\leqslant x\leqslant\pi, \\
u(0,t)=0,\quad u(\pi,t)=0, & t\geqslant 0;
\end{cases}$
\item $\begin{cases}
u_{tt}-a^2u_{xx}=0, & 0<x<l,\ t>0, \\
u(x,0)=x(x-2l),\quad u_t(x,0)=0, & 0\leqslant x\leqslant l, \\
u(0,t)=0,\quad u_x(l,t)=0, & t\geqslant 0;
\end{cases}$
\item $\begin{cases}
u_{tt}-a^2u_{xx}=0, & 0<x<l,\ t>0, \\
u(x,0)=\cos\frac{\pi x}{l},\quad u_t(x,0)=0, & 0\leqslant x\leqslant l, \\
u_x(0,t)=0,\quad u_x(l,t)=0, & t\geqslant 0;
\end{cases}$
\item $\begin{cases}
u_{tt}-a^2u_{xx}=0, & 0<x<l,\ t>0, \\
u(x,0)=1,\quad u_t(x,0)=1, & 0\leqslant x\leqslant l, \\
u_x(0,t)=0,\quad u_x(l,t)=0, & t\geqslant 0;
\end{cases}$
\item $\begin{cases}
u_{tt}-u_{xx}=0, & 0<x<\pi,\ t>0, \\
u(x,0)=0,\quad u_t(x,0)=0, & 0\leqslant x\leqslant\pi, \\
u_x(0,t)=A_1t,\quad u_x(\pi,t)=A_2t, & t\geqslant 0;
\end{cases}$
\item $\begin{cases}
u_{tt}-a^2u_{xx}=0, & 0<x<l,\ t>0, \\
u(x,0)=0,\quad u_t(x,0)=\cos\frac{\pi}{l}x, & 0\leqslant x\leqslant l, \\
u_x|_{x=0}=0,\quad (u_x+hu)|_{x=l}=0, & t\geqslant 0.
\end{cases}$
\end{enumerate}
\end{example}

\begin{example}
设 $a$ 为正常数，$A,\omega,A_1$ 为常数。用分离变量法求解下列混合问题：
\begin{enumerate}[(1)]
\item $\begin{cases}
u_{tt}-a^2u_{xx}=\sin\frac{\pi}{l}x, & 0<x<l,\ t>0, \\
u(x,0)=0,\quad u_t(x,0)=0, & 0\leqslant x\leqslant l, \\
u(0,t)=0,\quad u(l,t)=0, & t\geqslant 0;
\end{cases}$
\item $\begin{cases}
u_{tt}-a^2u_{xx}+2u_t=x(l-x), & 0<x<l,\ t>0, \\
u(x,0)=0,\quad u_t(x,0)=0, & 0\leqslant x\leqslant l, \\
u(0,t)=0,\quad u(l,t)=0, & t\geqslant 0;
\end{cases}$
\item $\begin{cases}
u_{tt}-a^2u_{xx}=0, & 0<x<l,\ t>0, \\
u(x,0)=\frac{Ax^2}{l^2},\quad u_t(x,0)=0, & 0\leqslant x\leqslant l, \\
u_x(0,t)=0,\quad u(l,t)=A, & t\geqslant 0;
\end{cases}$
\item $\begin{cases}
u_{tt}-a^2u_{xx}=0, & 0<x<l,\ t>0, \\
u(x,0)=1,\quad u_t(x,0)=0, & 0\leqslant x\leqslant l, \\
u_x(0,t)=A\sin\omega t,\quad u(l,t)=1, & t\geqslant 0,
\end{cases}$
分别讨论产生共振和不产生共振两种情形；
\item $\begin{cases}
u_{tt}-u_{xx}=0, & 0<x<l,\ t>0, \\
u(x,0)=0,\quad u_t(x,0)=0, & 0\leqslant x\leqslant l, \\
u_x|_{x=0}=A_1t,\quad (u_x+u)|_{x=l}=0, & t\geqslant 0.
\end{cases}$
\end{enumerate}
\end{example}

\begin{example}
求解混合问题
\begin{align*}
\begin{cases}
u_{tt}-u_{xx}-\frac{2}{x}u_x=0, & 0<x<1,\ t>0, \\
u(x,0)=0,\quad u_t(x,0)=\cos\frac{\pi}{2}x, & 0\leqslant x\leqslant 1, \\
u(0,t)=u(1,t)=0, & t\geqslant 0.
\end{cases}
\end{align*}
（提示：函数 $v(x,t)=xu(x,t)$ 满足 $v_{tt}-v_{xx}=0$）
\end{example}

\begin{example}
用分离变量法求解混合问题
\begin{align*}
u_{tt}-(u_{xx}+u_{yy})=0,\quad 0<x<1,\ 0<y<1,\ t>0,
\end{align*}
\begin{align*}
u|_{t=0}=x(1-x),\quad 0\leqslant x\leqslant 1,\ 0\leqslant y\leqslant 1,
\end{align*}
\begin{align*}
u_t|_{t=0}=y(1-y),\quad 0\leqslant x\leqslant 1,\ 0\leqslant y\leqslant 1,
\end{align*}
\begin{align*}
u|_{x=0}=u|_{x=1}=0,\quad 0\leqslant y\leqslant 1,\ t\geqslant 0,
\end{align*}
\begin{align*}
u_y|_{y=0}=u_y|_{y=1}=0,\quad 0\leqslant x\leqslant 1,\ t\geqslant 0.
\end{align*}
\end{example}

\begin{example}
设 $u(x,t)$ 满足混合问题
\begin{align*}
\begin{cases}
u_{tt}-a^2u_{xx}=f(x,t), & 0<x<l,\ t>0, \\
u(x,0)=\varphi(x),\quad u_t(x,0)=\psi(x), & 0\leqslant x\leqslant l, \\
(-u_x+\alpha u)|_{x=0}=g_1(t), & t\geqslant 0, \\
(u_x+\beta u)|_{x=l}=g_2(t), & t\geqslant 0,
\end{cases}
\end{align*}
其中 $a$ 为正常数。试引进辅助函数，将边值条件齐次化。设
\begin{enumerate}[(1)]
\item $\alpha>0,\ \beta>0$；
\item $\alpha=\beta=0$。
\end{enumerate}
\end{example}

在以下问题中记 $Q_T=(0,l)\times(0,T)$。

\begin{example}
设 $u(x,t)\in C^1(\overline{Q}_T)\cap C^2(\overline{Q}_T)$ 满足混合问题
\begin{align*}
\begin{cases}
u_{tt}-a^2u_{xx}=f(x,t), & (x,t)\in Q_T, \\
u(x,0)=\varphi(x),\quad u_t(x,0)=\psi(x), & 0\leqslant x\leqslant l, \\
u_x(0,t)=0,\quad u_x(l,t)=0, & 0\leqslant t\leqslant T,
\end{cases}
\end{align*}
其中 $a$ 为正常数。试推导能量不等式。
\end{example}

\begin{example}
设 $u(x,t)\in C^1(\overline{Q}_T)\cap C^2(\overline{Q}_T)$ 满足混合问题
\begin{align*}
\begin{cases}
u_{tt}-a^2u_{xx}=f(x,t), & (x,t)\in Q_T, \\
u(x,0)=\varphi(x),\quad u_t(x,0)=\psi(x), & 0\leqslant x\leqslant l, \\
(-u_x+\alpha u)|_{x=0}=0,\quad (u_x+\beta u)|_{x=l}=0, & 0\leqslant t\leqslant T,
\end{cases}
\end{align*}
其中 $a,\alpha,\beta>0$ 为常数。试推导能量不等式。
\end{example}

\begin{example}
设 $u(x,t)\in C^1(\overline{Q}_T)\cap C^2(\overline{Q}_T)$ 满足混合问题
\begin{align*}
\begin{cases}
u_{tt}-a^2u_{xx}=f(x,t), & (x,t)\in Q_T, \\
u(x,0)=\varphi(x),\quad u_t(x,0)=\psi(x), & 0\leqslant x\leqslant l, \\
u_x|_{x=0}=0,\quad (u_x+u)|_{x=l}=0, & 0\leqslant t\leqslant T,
\end{cases}
\end{align*}
其中 $a$ 为正常数。试推导能量不等式。
\end{example}

\begin{example}
考虑定解问题
\begin{align*}
\begin{cases}
u_{tt}-u_{xx}=f(x,t), & (x,t)\in Q_T, \\
u(x,0)=\varphi(x),\quad u_t(x,0)=\psi(x), & 0\leqslant x\leqslant l, \\
u(0,t)=u(l,t)=0, & 0\leqslant t\leqslant T.
\end{cases}
\end{align*}
试问：对 $\varphi,\psi,f$ 提什么条件才能保证由分离变量法所得的解是古典解？试证明之。
\end{example}

\begin{example}
试引入定解问题
\begin{align*}
\begin{cases}
u_{tt}-u_{xx}=f(x,t), & (x,t)\in Q_T, \\
u(x,0)=\varphi(x),\quad u_t(x,0)=\psi(x), & 0\leqslant x\leqslant l, \\
u(0,t)=u(l,t)=0, & 0\leqslant t\leqslant T
\end{cases}
\end{align*}
的广义解的定义。试问：对 $\varphi,\psi,f$ 提什么条件才能保证上述定解问题的广义解的存在性唯一性？试证明之。
\end{example}

\begin{example}
设 $u(x,t)\in C(\overline{Q}_T)$ 是混合问题 \eqref{eq:wave-gen-mixed} 的广义解。证明：
\begin{enumerate}[(1)]
\item $u(x,0)=\varphi(x),\quad x\in[0,l]$；
\item 如果 $u(x,t)\in C^1(\overline{Q}_T)$，则 $u_t(x,0)=\psi(x),\quad x\in[0,l]$。
\end{enumerate}
\end{example}




\end{document}